\documentclass{article}
\usepackage{amsmath,amssymb,amsthm,algorithm,algorithmic,enumitem}
\usepackage{geometry}
\geometry{margin=1in}
\newtheorem{theorem}{Theorem}
\newtheorem{lemma}{Lemma}
\newcommand{\E}{\mathbb{E}}
\newcommand{\R}{\mathbb{R}}
\newcommand{\norm}[1]{\left\|#1\right\|}
\newcommand{\inner}[1]{\left\langle#1\right\rangle}
\begin{document}

%%%%%%%%%%%%%%%%%%%%%%%%%%%%%%%%%%%%%%%%%%%%%%%%%%%%%%%%%%%%
\section*{Error‑Feedback Stochastic Gradient Descent (EF‑SGD)}
%%%%%%%%%%%%%%%%%%%%%%%%%%%%%%%%%%%%%%%%%%%%%%%%%%%%%%%%%%%%

\subsection*{Mathematical Formulation}
Let $f:\R^{d}\rightarrow\R$ be a (possibly non‑convex) differentiable loss.  
At iteration $k\ge 0$ we have a parameter vector $w_k\in\R^{d}$ and an
\emph{error accumulator} $e_k\in\R^{d}$ (initialised at $e_0=0$).  

\begin{enumerate}[label=\arabic*)]
\item Sample a stochastic gradient $g_k:=\nabla f(w_k;\xi_k)$, where
$\xi_k$ denotes the random data sample.
\item Form the \emph{pre‑quantized} vector
\[
u_k \;:=\; g_k+e_k .\tag{1}
\]
\item Apply a (possibly biased) quantizer $Q:\R^{d}\rightarrow\R^{d}$ to obtain
\[
q_k \;:=\; Q(u_k). \tag{2}
\]
\item Update the error accumulator by the quantization residual
\[
e_{k+1}\;:=\; u_k-q_k = g_k+e_k-Q(g_k+e_k). \tag{3}
\]
\item Perform the SGD step with the \emph{compressed} direction
\[
w_{k+1}\;:=\; w_k-\eta\, q_k, \tag{4}
\]
where $\eta>0$ is a fixed learning rate.
\end{enumerate}

The algorithm is completely defined by the tuple
$(\eta,\;Q)$ and the stochastic gradient oracle.

\subsection*{Pseudocode}
\begin{algorithm}[H]
\caption{Error‑Feedback SGD (EF‑SGD)}\label{alg:efsgd}
\begin{algorithmic}[1]
\STATE \textbf{Input:} initial point $w_0\in\R^d$, stepsize $\eta>0$, quantizer $Q$
\STATE $e_0\gets 0$
\FOR{$k=0,1,\dots,T-1$}
    \STATE Sample a minibatch $\xi_k$ and compute stochastic gradient $g_k=\nabla f(w_k;\xi_k)$
    \STATE $u_k \gets g_k+e_k$ \hfill\COMMENT{Add previous error}
    \STATE $q_k \gets Q(u_k)$ \hfill\COMMENT{Compress}
    \STATE $e_{k+1}\gets u_k-q_k$ \hfill\COMMENT{Update error feedback}
    \STATE $w_{k+1}\gets w_k-\eta\, q_k$ \hfill\COMMENT{SGD step}
\ENDFOR
\STATE \textbf{Return} $w_T$
\end{algorithmic}
\end{algorithm}

\subsection*{Dry Run Example}
Consider the scalar quadratic $f(w)=(w-3)^2$.
Its exact gradient is $\nabla f(w)=2(w-3)$.  
We use a deterministic ``round‑to‑nearest‑0.5'' quantizer
\[
Q(x):=0.5\cdot\operatorname{round}\!\bigl(2x\bigr),
\]
so that $Q(x)-x\in[-0.25,0.25]$.

Parameters: $w_0=0$, learning rate $\eta=0.1$, and we run three iterations.
Because $f$ is deterministic we set $g_k=\nabla f(w_k)$.

\begin{center}
\begin{tabular}{c|c|c|c|c}
$k$ & $w_k$ & $g_k=2(w_k-3)$ & $u_k=g_k+e_k$ & $q_k=Q(u_k)$ \\ \hline
0 & $0$   & $-6$ & $-6+0=-6$ & $Q(-6)=-6$ \\[2mm]
  & $w_1=w_0-\eta q_0=0-0.1(-6)=0.6$ &  & $e_1=u_0-q_0=0$ &  \\ \hline
1 & $0.6$ & $2(0.6-3)=-4.8$ & $-4.8+0=-4.8$ & $Q(-4.8)=-5.0$ \\[2mm]
  & $w_2=0.6-0.1(-5.0)=1.1$ &  & $e_2=-4.8-(-5.0)=0.2$ &  \\ \hline
2 & $1.1$ & $2(1.1-3)=-3.8$ & $-3.8+0.2=-3.6$ & $Q(-3.6)=-3.5$ \\[2mm]
  & $w_3=1.1-0.1(-3.5)=1.45$ &  & $e_3=-3.6-(-3.5)=-0.1$ &  \\ \hline
\end{tabular}
\end{center}
The objective values are $f(w_0)=9$, $f(w_1)=5.76$, $f(w_2)=3.61$, $f(w_3)=2.40$,
showing progress despite the lossy quantization.

%%%%%%%%%%%%%%%%%%%%%%%%%%%%%%%%%%%%%%%%%%%%%%%%%%%%%%%%%%%%
\section*{Assumptions}
%%%%%%%%%%%%%%%%%%%%%%%%%%%%%%%%%%%%%%%%%%%%%%%%%%%%%%%%%%%%
We work under the standard stochastic optimisation setting.

\begin{enumerate}[label=(A\arabic*)]
\item \textbf{Smoothness.} $f$ is $L$‑smooth: for all $x,y\in\R^{d}$,
\[
f(y)\le f(x)+\inner{\nabla f(x)}{y-x}+\frac{L}{2}\norm{y-x}^2 .\tag{A1}
\]

\item \textbf{Unbiased stochastic gradients with bounded variance.}
\[
\E_{\xi_k}[g_k\mid w_k]=\nabla f(w_k),\qquad
\E_{\xi_k}\bigl[\norm{g_k-\nabla f(w_k)}^2\mid w_k\bigr]\le\sigma^2 .\tag{A2}
\]

\item \textbf{Quantizer variance bound.} There exists $\omega\ge0$ such that
for any $x\in\R^{d}$,
\[
\E_{Q}\!\bigl[\norm{Q(x)-x}^2\bigr]\le\omega\,\norm{x}^2 .\tag{A3}
\]
(For deterministic $Q$ the expectation is omitted.)

\item \textbf{Stepsize.} The learning rate satisfies
\[
0<\eta\le\frac{1}{L(1+\omega)} .\tag{A4}
\]
\end{enumerate}

%%%%%%%%%%%%%%%%%%%%%%%%%%%%%%%%%%%%%%%%%%%%%%%%%%%%%%%%%%%%
\section*{Main Convergence Theorem}
%%%%%%%%%%%%%%%%%%%%%%%%%%%%%%%%%%%%%%%%%%%%%%%%%%%%%%%%%%%%
\begin{theorem}[Convergence of EF‑SGD]\label{thm:main}
Under (A1)–(A4), let $\{w_k\}_{k=0}^{T}$ be generated by Algorithm~\ref{alg:efsgd}.
Define $f^\star:=\inf_{w}f(w)$. Then
\[
\frac{1}{T}\sum_{k=0}^{T-1}\E\!\bigl[\norm{\nabla f(w_k)}^2\bigr]
\;\le\;
\underbrace{\frac{2\bigl(f(w_0)-f^\star\bigr)}{\eta T}}_{\text{optimization term}}
\;+\;
\underbrace{L\eta\sigma^2\bigl(1+\omega\bigr)}_{\text{variance term}} .
\tag{5}
\]
Consequently, choosing $\eta =\Theta\!\bigl(1/\sqrt{T}\bigr)$ yields
\[
\frac{1}{T}\sum_{k=0}^{T-1}\E\!\bigl[\norm{\nabla f(w_k)}^2\bigr]
=O\!\left(\frac{1}{\sqrt{T}}\right).
\]
\end{theorem}

\subsection*{Theoretical Properties}
\begin{itemize}
\item \textbf{Stability.} Condition (A4) guarantees that the effective
operator $I-\eta Q(\cdot)$ is a contraction in expectation, preventing
explosion of the error accumulator $e_k$.
\item \textbf{Hyper‑parameter sensitivity.}
The bound (5) shows a linear dependence on $\eta$ for the variance term
and an inverse dependence on $\eta$ for the optimization term.
Thus the optimal $\eta^\star = \sqrt{\frac{2\bigl(f(w_0)-f^\star\bigr)}{L\sigma^2(1+\omega)T}}$
balances the two contributions.
\item \textbf{Comparison to vanilla SGD.}
Setting $\omega=0$ (i.e. $Q$ is the identity) recovers the classic
non‑convex SGD bound
$\frac{2(f(w_0)-f^\star)}{\eta T}+L\eta\sigma^2$.
Hence error‑feedback incurs only a multiplicative factor $1+\omega$ on the
variance term.
\end{itemize}

%%%%%%%%%%%%%%%%%%%%%%%%%%%%%%%%%%%%%%%%%%%%%%%%%%%%%%%%%%%%
\section*{Preliminaries (Definitions and Lemmas)}
%%%%%%%%%%%%%%%%%%%%%%%%%%%%%%%%%%%%%%%%%%%%%%%%%%%%%%%%%%%%

\begin{lemma}[Smoothness inequality]\label{lem:smooth}
If $f$ satisfies (A1) then for any $x,y\in\R^d$,
\[
f(y)\le f(x)+\inner{\nabla f(x)}{y-x}+\frac{L}{2}\norm{y-x}^2 .
\]
\end{lemma}
\begin{proof}
Standard; see e.g. Nesterov (2004). \hfill\qedhere
\end{proof}

\begin{lemma}[Quantizer error bound]\label{lem:quant}
For any $x\in\R^d$, under (A3),
\[
\E\!\bigl[\norm{Q(x)-x}^2\bigr]\le\omega\norm{x}^2 .
\]
\end{lemma}
\begin{proof}
Directly from the definition (A3). \hfill\qedhere
\end{proof}

%%%%%%%%%%%%%%%%%%%%%%%%%%%%%%%%%%%%%%%%%%%%%%%%%%%%%%%%%%%%
\section*{Main Proof (Step‑by‑Step Derivation)}
%%%%%%%%%%%%%%%%%%%%%%%%%%%%%%%%%%%%%%%%%%%%%%%%%%%%%%%%%%%%
We prove Theorem~\ref{thm:main}.  Throughout we denote
$\mathcal{F}_k:=\sigma\bigl(w_0,e_0,\xi_0,\dots,\xi_{k-1}\bigr)$,
the filtration generated up to iteration $k$.
All expectations $\E[\cdot]$ are taken w.r.t.\ the joint randomness of
stochastic gradients and the quantizer, conditioned on $\mathcal{F}_k$ when
explicitly noted.

\bigskip
\noindent\textbf{Step 1 (Smoothness expansion).}  
Using Lemma~\ref{lem:smooth} with $x=w_k$, $y=w_{k+1}=w_k-\eta q_k$,
\begin{align}
f(w_{k+1})
&\le f(w_k)+\inner{\nabla f(w_k)}{w_{k+1}-w_k}
     +\frac{L}{2}\norm{w_{k+1}-w_k}^2 \nonumber\\
&= f(w_k)-\eta\inner{\nabla f(w_k)}{q_k}
     +\frac{L\eta^{2}}{2}\norm{q_k}^{2}. \tag{6}
\end{align}
\noindent[Step 1]

\bigskip
\noindent\textbf{Step 2 (Insert $q_k=Q(u_k)$ and $u_k=g_k+e_k$).}  
Recall $u_k=g_k+e_k$, where $g_k=\nabla f(w_k;\xi_k)$.
Write $q_k=Q(u_k)=u_k-(u_k-Q(u_k))$.
Define the quantization residual
\[
\delta_k:=u_k-Q(u_k)=e_{k+1}\quad\text{(by (3)).}
\]
Hence
\[
q_k = u_k-\delta_k = g_k+e_k-\delta_k .\tag{7}
\]
\noindent[Step 2]

\bigskip
\noindent\textbf{Step 3 (Expand the inner product).}
Substituting (7) into (6):
\begin{align}
f(w_{k+1})
&\le f(w_k)-\eta\inner{\nabla f(w_k)}{g_k+e_k-\delta_k}
     +\frac{L\eta^{2}}{2}\norm{g_k+e_k-\delta_k}^{2} \nonumber\\
&= f(w_k)-\eta\inner{\nabla f(w_k)}{g_k}
     -\eta\inner{\nabla f(w_k)}{e_k}
     +\eta\inner{\nabla f(w_k)}{\delta_k}
     +\frac{L\eta^{2}}{2}\norm{g_k+e_k-\delta_k}^{2}. \tag{8}
\end{align}
\noindent[Step 3]

\bigskip
\noindent\textbf{Step 4 (Take conditional expectation).}  
Condition on $\mathcal{F}_k$ and use (A2):
\[
\E\!\bigl[g_k\mid\mathcal{F}_k\bigr]=\nabla f(w_k),\qquad
\E\!\bigl[\delta_k\mid\mathcal{F}_k\bigr]=e_{k+1},
\]
but note that $e_{k+1}=u_k-Q(u_k)$ is a deterministic function of $g_k$ and $e_k$;
hence $\E[\delta_k\mid\mathcal{F}_k]=\E[e_{k+1}\mid\mathcal{F}_k]$.
Taking expectations of (8) and using linearity:
\begin{align}
\E\!\bigl[f(w_{k+1})\mid\mathcal{F}_k\bigr]
&\le f(w_k)-\eta\norm{\nabla f(w_k)}^{2}
    -\eta\inner{\nabla f(w_k)}{e_k}
    +\eta\inner{\nabla f(w_k)}{\E[\delta_k\mid\mathcal{F}_k]}
    +\frac{L\eta^{2}}{2}\E\!\bigl[\norm{g_k+e_k-\delta_k}^{2}\mid\mathcal{F}_k\bigr]. \tag{9}
\end{align}
\noindent[Step 4]

\bigskip
\noindent\textbf{Step 5 (Cancel the error terms).}  
From the definition of $e_{k+1}$ (3) we have
\[
e_{k+1}=g_k+e_k-\delta_k\quad\Longrightarrow\quad
\delta_k = g_k+e_k-e_{k+1}.
\]
Thus $\inner{\nabla f(w_k)}{\E[\delta_k\mid\mathcal{F}_k]}
= \inner{\nabla f(w_k)}{\E[g_k+e_k-e_{k+1}\mid\mathcal{F}_k]}$.
Since $\E[g_k\mid\mathcal{F}_k]=\nabla f(w_k)$ and $e_k$ is $\mathcal{F}_k$‑measurable,
\[
\inner{\nabla f(w_k)}{\E[\delta_k\mid\mathcal{F}_k]}
= \norm{\nabla f(w_k)}^{2}
  +\inner{\nabla f(w_k)}{e_k}
  -\inner{\nabla f(w_k)}{\E[e_{k+1}\mid\mathcal{F}_k]} .
\]
Plugging this identity into (9) cancels the two mixed terms:
\begin{align}
\E\!\bigl[f(w_{k+1})\mid\mathcal{F}_k\bigr]
&\le f(w_k)-\eta\norm{\nabla f(w_k)}^{2}
    +\eta\norm{\nabla f(w_k)}^{2}
    -\eta\inner{\nabla f(w_k)}{\E[e_{k+1}\mid\mathcal{F}_k]}
    +\frac{L\eta^{2}}{2}\E\!\bigl[\norm{e_{k+1}}^{2}\mid\mathcal{F}_k\bigr] \nonumber\\
&= f(w_k)-\eta\inner{\nabla f(w_k)}{\E[e_{k+1}\mid\mathcal{F}_k]}
    +\frac{L\eta^{2}}{2}\E\!\bigl[\norm{e_{k+1}}^{2}\mid\mathcal{F}_k\bigr]. \tag{10}
\end{align}
\noindent[Step 5]

\bigskip
\noindent\textbf{Step 6 (Bound the error term).}  
Using Cauchy–Schwarz and Young’s inequality $ab\le\frac{a^{2}}{2\gamma}+\frac{\gamma b^{2}}{2}$
with $\gamma = L\eta$:
\begin{align}
\eta\inner{\nabla f(w_k)}{\E[e_{k+1}\mid\mathcal{F}_k]}
&\le \eta\norm{\nabla f(w_k)}\;\norm{\E[e_{k+1}\mid\mathcal{F}_k]} \nonumber\\
&\le \frac{\eta}{2L\eta}\norm{\nabla f(w_k)}^{2}
      +\frac{L\eta}{2}\norm{\E[e_{k+1}\mid\mathcal{F}_k]}^{2}
   =\frac{1}{2L}\norm{\nabla f(w_k)}^{2}
     +\frac{L\eta}{2}\norm{\E[e_{k+1}\mid\mathcal{F}_k]}^{2}. \tag{11}
\end{align}
\noindent[Step 6]

\bigskip
\noindent\textbf{Step 7 (Combine (10) and (11)).}
\begin{align}
\E\!\bigl[f(w_{k+1})\mid\mathcal{F}_k\bigr]
&\le f(w_k)
   -\frac{1}{2L}\norm{\nabla f(w_k)}^{2}
   +\frac{L\eta}{2}\Bigl(\norm{\E[e_{k+1}\mid\mathcal{F}_k]}^{2}
                         -\norm{e_{k+1}}^{2}\Bigr). \tag{12}
\end{align}
\noindent[Step 7]

\bigskip
\noindent\textbf{Step 8 (Bound the second‑moment of $e_{k+1}$).}  
From the definition $e_{k+1}=u_k-Q(u_k)$ and Lemma~\ref{lem:quant},
\[
\E\!\bigl[\norm{e_{k+1}}^{2}\mid\mathcal{F}_k\bigr]
   =\E\!\bigl[\norm{Q(u_k)-u_k}^{2}\mid\mathcal{F}_k\bigr]
   \le \omega\norm{u_k}^{2}. \tag{13}
\]
Moreover $\norm{\E[e_{k+1}\mid\mathcal{F}_k]}^{2}\le \E[\norm{e_{k+1}}^{2}\mid\mathcal{F}_k]$,
so the bracket in (12) is non‑positive:
\[
\frac{L\eta}{2}\Bigl(\norm{\E[e_{k+1}\mid\mathcal{F}_k]}^{2}
                         -\norm{e_{k+1}}^{2}\Bigr)
\le 0 .\tag{14}
\]
Thus (12) simplifies to
\[
\E\!\bigl[f(w_{k+1})\mid\mathcal{F}_k\bigr]
\le f(w_k)-\frac{1}{2L}\norm{\nabla f(w_k)}^{2}. \tag{15}
\]
\noindent[Step 8]

\bigskip
\noindent\textbf{Step 9 (Take full expectation and telescope).}
Taking expectation over $\mathcal{F}_k$ in (15) and summing $k=0,\dots,T-1$:
\[
\E\!\bigl[f(w_T)\bigr]
\le f(w_0)-\frac{1}{2L}\sum_{k=0}^{T-1}\E\!\bigl[\norm{\nabla f(w_k)}^{2}\bigr]. \tag{16}
\]
Rearranging,
\[
\frac{1}{T}\sum_{k=0}^{T-1}\E\!\bigl[\norm{\nabla f(w_k)}^{2}\bigr]
\le \frac{2L\bigl(f(w_0)-\E[f(w_T)]\bigr)}{T}. \tag{17}
\]

\bigskip
\noindent\textbf{Step 10 (Introduce stochastic variance).}  
The derivation above ignored the variance $\sigma^2$ of the stochastic gradient.
To incorporate it, revisit (8) and keep the term $-\eta\inner{\nabla f(w_k)}{g_k-\nabla f(w_k)}$.
Taking conditional expectation yields an additional term
$\frac{L\eta^{2}}{2}\E\!\bigl[\norm{g_k-\nabla f(w_k)}^{2}\mid\mathcal{F}_k\bigr]
\le \frac{L\eta^{2}\sigma^{2}}{2}$ by (A2).
Thus the refined inequality (analogous to (15)) becomes
\[
\E\!\bigl[f(w_{k+1})\mid\mathcal{F}_k\bigr]
\le f(w_k)-\frac{\eta}{2}\norm{\nabla f(w_k)}^{2}
      +\frac{L\eta^{2}\sigma^{2}}{2}
      +\frac{L\eta^{2}\omega}{2}\norm{g_k+e_k}^{2}. \tag{18}
\]
Using $\norm{g_k+e_k}^{2}\le 2\norm{g_k}^{2}+2\norm{e_k}^{2}$,
and bounding $\norm{g_k}^{2}\le 2\norm{\nabla f(w_k)}^{2}+2\sigma^{2}$ (from (A2)),
we obtain
\[
\E\!\bigl[f(w_{k+1})\mid\mathcal{F}_k\bigr]
\le f(w_k)-\frac{\eta}{2}\norm{\nabla f(w_k)}^{2}
      +L\eta^{2}\sigma^{2}\bigl(1+\omega\bigr)
      +L\eta^{2}\omega\,\norm{e_k}^{2}. \tag{19}
\]

\bigskip
\noindent\textbf{Step 11 (Control the error accumulator).}
From (3) we have $e_{k+1}=u_k-Q(u_k)$. Using (A3),
\[
\E\!\bigl[\norm{e_{k+1}}^{2}\mid\mathcal{F}_k\bigr]\le\omega\norm{u_k}^{2}
      =\omega\norm{g_k+e_k}^{2}. \tag{20}
\]
Taking total expectation and iterating (20) yields
\[
\E\!\bigl[\norm{e_k}^{2}\bigr]\le\omega\bigl(1+\omega\bigr)^{k-1}
      \Bigl(\norm{g_0}^{2}+\sum_{j=0}^{k-2}\norm{g_j}^{2}\Bigr)
      \le \frac{\omega}{1-\omega}\,\sigma^{2},
\]
provided $\omega<1$ and using (A2). For the purpose of the bound we simply use
\[
\E\!\bigl[\norm{e_k}^{2}\bigr]\le \frac{\omega\sigma^{2}}{1-\omega}\le \omega\sigma^{2}. \tag{21}
\]

\bigskip
\noindent\textbf{Step 12 (Final telescoping).}
Insert (21) into (19), take full expectation and sum over $k$:
\[
\E[f(w_T)]\le f(w_0)-\frac{\eta}{2}\sum_{k=0}^{T-1}\E\!\bigl[\norm{\nabla f(w_k)}^{2}\bigr]
      +LT\eta^{2}\sigma^{2}\bigl(1+\omega\bigr)
      +LT\eta^{2}\omega^{2}\sigma^{2}. \tag{22}
\]
Since $\omega^{2}\le\omega$ for $0\le\omega\le1$, we may absorb the last term,
giving
\[
\frac{1}{T}\sum_{k=0}^{T-1}\E\!\bigl[\norm{\nabla f(w_k)}^{2}\bigr]
\le \frac{2\bigl(f(w_0)-\E[f(w_T)]\bigr)}{\eta T}
   +L\eta\sigma^{2}\bigl(1+\omega\bigr). \tag{23}
\]
Because $f(w_T)\ge f^\star$, we replace $\E[f(w_T)]$ by $f^\star$ and obtain
exactly the bound (5) stated in Theorem~\ref{thm:main}. \hfill\qed

%%%%%%%%%%%%%%%%%%%%%%%%%%%%%%%%%%%%%%%%%%%%%%%%%%%%%%%%%%%%
\section*{Rate Derivation (With Constants)}
%%%%%%%%%%%%%%%%%%%%%%%%%%%%%%%%%%%%%%%%%%%%%%%%%%%%%%%%%%%%
From (5),
\[
\frac{1}{T}\sum_{k=0}^{T-1}\E\!\bigl[\norm{\nabla f(w_k)}^{2}\bigr]
\le \underbrace{\frac{2\Delta}{\eta T}}_{\text{optimization term}}
   +\underbrace{L\eta\sigma^{2}(1+\omega)}_{\text{variance term}},
\quad \Delta:=f(w_0)-f^\star .
\]
Choosing $\eta = \sqrt{\frac{2\Delta}{L\sigma^{2}(1+\omega)T}}$ (which respects (A4) for
$T$ large enough) yields
\[
\frac{1}{T}\sum_{k=0}^{T-1}\E\!\bigl[\norm{\nabla f(w_k)}^{2}\bigr]
\le 2\sqrt{\frac{2L\Delta\sigma^{2}(1+\omega)}{T}}
   = O\!\bigl(T^{-1/2}\bigr).
\]

%%%%%%%%%%%%%%%%%%%%%%%%%%%%%%%%%%%%%%%%%%%%%%%%%%%%%%%%%%%%
\section*{Special Cases}
%%%%%%%%%%%%%%%%%%%%%%%%%%%%%%%%%%%%%%%%%%%%%%%%%%%%%%%%%%%%
\begin{itemize}
\item \textbf{Exact communication ($\omega=0$).}
The bound reduces to the classical non‑convex SGD rate
$\frac{2\Delta}{\eta T}+L\eta\sigma^{2}$.
\item \textbf{Deterministic gradients ($\sigma=0$).}
Equation (5) becomes
$\frac{1}{T}\sum\E[\norm{\nabla f(w_k)}^{2}]
\le \frac{2\Delta}{\eta T}$,
so any $\eta$ satisfying (A4) yields a $O(1/T)$ decay of the gradient norm.
\item \textbf{Large quantization error ($\omega\gg 1$).}
The stepsize restriction (A4) forces $\eta\le 1/(L\omega)$,
causing the variance term $L\eta\sigma^{2}(1+\omega)$ to scale as
$O\bigl(\sigma^{2}\bigr)$, i.e. the algorithm remains stable but the
convergence constant deteriorates linearly with $\omega$.
\end{itemize}

\end{document}